\clearpage  % Appendix F opens a page, so the field table and the example row print beside their citation
\section{Dataset Fields}
\label{app:dataset-schema}

We release the audited benchmark after review as a \texttt{datasets}
dataset~\citep{lhoest2021datasets} with two splits and no \texttt{train}
split. The \texttt{test} split, marked \texttt{split="eval"}, is the 100-paper
frozen benchmark, and the \texttt{validation} split, marked \texttt{split="dev"},
is the 14-paper development set (Table~\ref{tab:splits}), and no arXiv ID
appears in both. Every row holds the 25 top-level fields and four nested objects of
Table~\ref{tab:schema}, fixed before any agent starts. Figure~\ref{lst:example-row} shows the SharpZO row of
Section~\ref{subsec:task} (availability score 0, Run tier, 0.07 H100-hours, band
\texttt{0-8}), long text fields truncated and the rest verbatim. Rows return
metadata and links only, with no full paper text.

\begin{table}[ht]
\centering
\small
\begin{tabular}{@{}lrrrrrr@{}}
\toprule
& \multicolumn{5}{c}{Evaluation split} & Development split \\
\cmidrule(lr){2-6} \cmidrule(lr){7-7}
Tier & Papers & 0--8 & 8--32 & 32--96 & Audited H100-h & Papers \\
\midrule
Run         & 34  & 27 & 7  & 0  & 142.71    & 5  \\
Retrain     & 33  & 16 & 13 & 4  & 380.91    & 4  \\
Reimplement & 33  & 16 & 9  & 8  & 630.18    & 5  \\
\midrule
All         & 100 & 59 & 29 & 12 & 1{,}153.8 & 14 \\
\bottomrule
\end{tabular}
\caption{Composition of the two splits by tier and compute band, with papers
binned by the audited H100-hour estimate.}
\label{tab:splits}
\end{table}

{\footnotesize
\begin{longtable}{@{}>{\raggedright\arraybackslash}p{4.95cm}>{\raggedright\arraybackslash}p{1.95cm}>{\raggedright\arraybackslash}p{5.65cm}@{}}
\toprule
Field & Type & Description \\
\midrule
\endfirsthead
\multicolumn{3}{@{}l}{\footnotesize\itshape Table~\ref{tab:schema}, continued.} \\[3pt]
\toprule
Field & Type & Description \\
\midrule
\endhead
\midrule
\endfoot
\bottomrule
\noalign{\vskip 6pt}
\caption{The fields of a dataset row.}
\label{tab:schema} \\
\endlastfoot
\multicolumn{3}{@{}>{\raggedright\arraybackslash}p{13.4cm}@{}}{\emph{Top-level fields}} \\
\texttt{custom\_id} & string & arXiv ID, the primary key. \\
\texttt{central\_claim} & string & The paper's central empirical claim in one sentence. \\
\texttt{claim\_evidence} & string & Quoted paper text grounding the claim. \\
\texttt{mre\_config} & string & Prose recipe for the cheapest configuration that tests the claim, withheld from the agent. \\
\texttt{agent\_task} & string & Classifier-written reproduction instruction, released with the row and shown to neither the agent nor the auditor. \\
\texttt{verified\_links} & object & Tool-confirmed artifact URLs. \\
\texttt{signals} & object & Per-artifact availability audit. \\
\texttt{match\_target} & object & The quantity a run must match. \\
\texttt{h100\_estimate} & object & Compute estimate and its basis. \\
\texttt{score} & int & Artifact-availability score from the four signals; shipped rows take 0 to 3. \\
\texttt{tier} & string & \texttt{Easy}, \texttt{Medium} or \texttt{Hard} in the released dataset, for the Run, Retrain and Reimplement tiers; derived from \texttt{score}. \\
\texttt{split} & string & \texttt{eval} (test split) or \texttt{dev} (validation split). \\
\texttt{selection\_band} & string & Compute band used for stratified selection. \\
\texttt{h100\_band} & string & Compute band, one of \texttt{0-8}, \texttt{8-32}, \texttt{32-96}. \\
\texttt{h100\_hours\_estimate} & float & Estimated H100-hours to run the match target. \\
\texttt{h100\_estimate\_basis} & string & Free-text justification for the estimate. \\
\texttt{audited\_h100\_hours} & float & Recomputed product where the arithmetic audit flagged the stated hours and the selection code accepted it, otherwise the stated estimate (Appendix~\ref{app:selection}). \\
\texttt{h100\_recomputed\_hours} & float or null & Hours recomputed from the estimate's arithmetic; null when an input is missing. \\
\texttt{h100\_arithmetic\_mismatch} & bool or null & Flag set when estimate arithmetic disagreed. \\
\texttt{h100\_hours\_adjudicated} & bool & Whether the recompute replaced the stated estimate. \\
\texttt{h100\_needs\_human\_review} & bool & Whether the estimate was flagged for review. \\
\texttt{paper\_kind} & string & Paper type. All rows are \texttt{empirical}. \\
\texttt{verification\_status} & string & Classifier verification state. All rows are \texttt{verified}. \\
\texttt{web\_verification} & string & Web-retrieval state. All rows are \texttt{available}. \\
\texttt{exit\_reason} & string & Classifier exit reason. All rows are \texttt{natural}. \\
\addlinespace
\multicolumn{3}{@{}>{\raggedright\arraybackslash}p{13.4cm}@{}}{\emph{\texttt{verified\_links} (lists of tool-confirmed URLs, empty when none located)}} \\
\texttt{paper\_or\_project} & list of strings & Paper or project-page URLs. \\
\texttt{code} & list of strings & Code repository URLs. \\
\texttt{dataset} & list of strings & Dataset URLs. \\
\texttt{weights} & list of strings & Model-weight URLs. \\
\addlinespace
\multicolumn{3}{@{}>{\raggedright\arraybackslash}p{13.4cm}@{}}{\emph{\texttt{signals} (one object per audited signal: \texttt{code\_available}, \texttt{dataset\_available}, \texttt{weights\_available}, \texttt{dataset\_is\_standard})}} \\*
\texttt{value} & bool & Whether the artifact is present. \\
\texttt{verification} & string & \texttt{tool\_verified}, \texttt{tool\_searched\_not\_found} or \texttt{not\_applicable}. \\
\texttt{evidence} & string & The tool observation supporting the value. \\
\addlinespace
\multicolumn{3}{@{}>{\raggedright\arraybackslash}p{13.4cm}@{}}{\emph{\texttt{match\_target} (the quantity a run must recover to count as a match)}} \\
\texttt{config} & string & Experimental configuration that produces the target value. \\
\texttt{metric} & string & The reported metric name. \\
\texttt{value} & string & The target value to match. \\
\texttt{scope} & string & The model, dataset or setting the value applies to. \\
\texttt{match\_bar\_kind} & string & \texttt{point\_estimate}, \texttt{direction}, \texttt{threshold}, \texttt{magnitude} or \texttt{range}. \\
\addlinespace
\multicolumn{3}{@{}>{\raggedright\arraybackslash}p{13.4cm}@{}}{\emph{\texttt{h100\_estimate} (the compute estimate that placed the paper in its band)}} \\
\texttt{hours} & float & Estimated H100-hours for the match target. \\
\texttt{basis\_kind} & string & \texttt{derived\_from\_config}, \texttt{paper\_reported}, \texttt{comparable\_experiment} or \texttt{compute\_unspecified}. \\
\texttt{gpu\_count} & int or null & GPU count the paper reports. \\
\texttt{gpu\_type} & string or null & GPU type the paper reports. \\
\texttt{wallclock\_hours} & float or null & Reported wall-clock hours, before H100 normalization. \\
\texttt{h100\_equivalent\_multiplier} & float or null & Factor converting the reported GPU to H100-equivalent hours. \\
\texttt{basis} & string & Free-text justification for the estimate. \\
\end{longtable}
}

\codefile[lst:example-row]{json}{One evaluation row (2506.20990)}{prompts/example_row.json}

\paragraph{Match targets and tolerance.} For 62 of the 100 evaluation papers
the match bar is a point estimate, for 21 a stated baseline the run must beat,
for 14 a reported threshold it must clear, for 1 the size of a reported delta,
and for 1 a reported range it must fall inside. A run matches a
point estimate when its measured value lies within the authors' stated
tolerance, or within 5\% relative when the paper states none
(Figure~\ref{lst:rubric-text}), and 58 of the 62 point estimates use the 5\%
default. The table uses the grades as submitted, before the corrections of Appendix~\ref{app:outside-agent} and Appendix~\ref{app:target-check}. To measure how much the default decides, we rescored every graded run
at relative tolerances of 1\%, 2\%, and 10\% (Table~\ref{tab:tolerance}). The
Reimplement tier has the lowest best-agent count at every setting, and the best agent's count moves
by at most 5 papers at Run, 3 at Retrain, and 1 at Reimplement. Tightening to
1\% moves 17 of the 73 runs reproduced as submitted to partial, and loosening to 10\% moves
up 9 partial runs whose only failed criterion was the number. For example,
DeepSeek-V4-Flash measured 90.6 on 2511.19808 against a frozen target of 97.3,
a gap of $6.7 / 97.3 = 6.9\%$, so that run is partial at 5\% and reproduced at
10\%. With an absolute tolerance of 1 or 2 points on percentage metrics, 70 or
72 of the 400 cells reproduce.
% src: results/e6/summary.json (claims: 62 point estimates, point_estimate_default_5pct 58, point_estimate_stated_margin 4; falls_by_setting rel1 17; number_only_partials 9; settings abs1, abs2); results/e6/README.md:13-24 (worked run; best-agent table)

\begin{table}[ht]
\centering
\small
\begin{tabular}{@{}lrrrr@{}}
\toprule
Relative tolerance & Run, best of 34 & Retrain, best of 33 & Reimplement, best of 33 & All agents, of 400 \\
\midrule
1\%           &  9 &  8 & 4 & 56 \\
2\%           & 12 &  9 & 5 & 65 \\
5\% (default) & 14 &  9 & 5 & 73 \\
10\%          & 14 & 12 & 5 & 82 \\
\bottomrule
\end{tabular}
\caption{Reproduced cells at four relative tolerances, for the best agent in
each tier and over all agents.}
\label{tab:tolerance}
\end{table}
% src: results/e6/README.md:19-24 (best agent per tier); results/e6/summary.json cells summed per setting (56, 65, 73, 82)
% moves with: E7 and the regrade decision (the 5% row becomes the post-correction counts); rerun rel1/rel2/rel10 if the counts of record change

\paragraph{The evaluation and development papers.} Tables~\ref{tab:eval100}
and~\ref{tab:dev14} list the 100 evaluation and 14 development papers with the
match target each is graded against, where Band is the compute band in
H100-hours, H100-h the audited estimate for that paper's match target, and Metric, Bar,
and Target value the pinned match target, with Bar written unset on the one
evaluation row with no bar kind. The
development split totals 26.89 audited H100-hours.

% Generated by paper_latex/tables/gen_eval100.py. Do not edit by hand.
% src: outputs/v6/app_rebuild/eval_100.jsonl
{\scriptsize
\begin{longtable}{@{}l@{\hspace{4pt}}l@{\hspace{4pt}}>{\centering\arraybackslash}p{0.9cm}@{\hspace{4pt}}>{\raggedleft\arraybackslash}p{1.1cm}@{\hspace{8pt}}p{3.2cm}@{\hspace{5pt}}p{1.6cm}@{\hspace{5pt}}p{3.3cm}@{}}
\toprule
arXiv & Tier & Band & H100-h & Metric & Bar & Target value \\
\midrule
\endfirsthead
\multicolumn{7}{@{}l}{\footnotesize\itshape Table~\ref{tab:eval100}, continued.} \\[3pt]
\toprule
arXiv & Tier & Band & H100-h & Metric & Bar & Target value \\
\midrule
\endhead
\midrule
\endfoot
\bottomrule
\noalign{\vskip 6pt}
\caption{The 100 evaluation papers, one row each, sorted by tier, then compute band, then arXiv ID. Metric names and target values longer than 28 characters are cut and closed with an ellipsis.}
\label{tab:eval100} \\
\endlastfoot
2502.05795 & Run & 0--8 & 1.00 & \raggedright Validation perplexity... & \raggedright point estimate & \raggedright 25.76 \tabularnewline
2502.06684 & Run & 0--8 & $<$0.01 & \raggedright Mean accuracy (\%) on 10... & \raggedright point estimate & \raggedright 77.0 +/- 1.8\% \tabularnewline
2503.18430 & Run & 0--8 & 0.64 & \raggedright Average Precision (AP) & \raggedright point estimate & \raggedright 46.3\% \tabularnewline
2503.23035 & Run & 0--8 & 0.25 & \raggedright PSNR (dB) & \raggedright point estimate & \raggedright 27.69 \tabularnewline
2504.20571 & Run & 0--8 & 0.16 & \raggedright MATH500 accuracy (\%) & \raggedright point estimate & \raggedright 73.6\% \tabularnewline
2505.10978 & Run & 0--8 & 2.00 & \raggedright Success rate (\%) & \raggedright point estimate & \raggedright 86.7\% \tabularnewline
2505.11483 & Run & 0--8 & 0.00 & \raggedright Peak RAM usage & \raggedright point estimate & \raggedright 8.56 kB \tabularnewline
2505.14766 & Run & 0--8 & 1.00 & \raggedright Mean Absolute Scaled Error... & \raggedright point estimate & \raggedright 0.617 \tabularnewline
2505.14827 & Run & 0--8 & 8.00 & \raggedright Accuracy (\%) & \raggedright direction & \raggedright 59.53\% \tabularnewline
2505.18456 & Run & 0--8 & 2.50 & \raggedright test perplexity (PPL, lower... & \raggedright point estimate & \raggedright 21.66 \tabularnewline
2505.18809 & Run & 0--8 & 2.00 & \raggedright end-to-end video generation... & \raggedright point estimate & \raggedright 1.76x \tabularnewline
2505.19154 & Run & 0--8 & 0.03 & \raggedright PSNR (dB) on rendered novel... & \raggedright point estimate & \raggedright 30.9 \tabularnewline
2505.19713 & Run & 0--8 & 0.64 & \raggedright Mean Chamfer Distance ($\times$10$^3$) & \raggedright point estimate & \raggedright 6.54 \tabularnewline
2505.23747 & Run & 0--8 & 2.56 & \raggedright VSI-Bench micro average... & \raggedright point estimate & \raggedright 51.13\% \tabularnewline
2505.24864 & Run & 0--8 & 8.00 & \raggedright Average pass@1 improvement... & \raggedright direction & \raggedright 15.7\% \tabularnewline
2505.24873 & Run & 0--8 & 1.00 & \raggedright SSIM (structural similarity) & \raggedright point estimate & \raggedright 0.9842 \tabularnewline
2506.02392 & Run & 0--8 & $<$0.01 & \raggedright optimality gap (\%) & \raggedright point estimate & \raggedright 3.25\% \tabularnewline
2506.10351 & Run & 0--8 & 1.00 & \raggedright Classification accuracy (\%) & \raggedright point estimate & \raggedright 93.12\% \tabularnewline
2506.12025 & Run & 0--8 & 0.50 & \raggedright Prediction speedup... & \raggedright direction & \raggedright 100x \tabularnewline
2506.18890 & Run & 0--8 & $<$0.01 & \raggedright PSNR (dB) & \raggedright point estimate & \raggedright 27.443 \tabularnewline
2506.20671 & Run & 0--8 & 0.16 & \raggedright PQ-All (\%) & \raggedright point estimate & \raggedright 6.30 \tabularnewline
2506.20990 & Run & 0--8 & 0.07 & \raggedright test accuracy (\%) & \raggedright point estimate & \raggedright 79.42\% \tabularnewline
2506.21724 & Run & 0--8 & 2.00 & \raggedright Overall accuracy (\%) on... & \raggedright point estimate & \raggedright 90.53\% \tabularnewline
2510.21311 & Run & 0--8 & 0.16 & \raggedright gIoU (generic Intersection... & \raggedright point estimate & \raggedright 55.1 \tabularnewline
2510.23574 & Run & 0--8 & 3.00 & \raggedright absolute relative error... & \raggedright threshold & \raggedright 7.5 \tabularnewline
2511.00090 & Run & 0--8 & 1.00 & \raggedright Latency speedup ratio & \raggedright point estimate & \raggedright 2.59x \tabularnewline
2511.16666 & Run & 0--8 & 2.00 & \raggedright Orientation accuracy at... & \raggedright point estimate & \raggedright 89.47\% \tabularnewline
2410.18164 & Run & 8--32 & 18.00 & \raggedright CC18 AUC / CTR23 R2 & \raggedright threshold & \raggedright AUC \textgreater{}= 0.92 (paper 0.933)... \tabularnewline
2412.11979 & Run & 8--32 & 16.00 & \raggedright Zipf power-law exponent... & \raggedright range & \raggedright approx 0.8-1.0 (observed... \tabularnewline
2504.12397 & Run & 8--32 & 10.00 & \raggedright Classification accuracy (\%) & \raggedright direction & \raggedright 79.30\% (LoRA) vs 81.94\%... \tabularnewline
2505.10475 & Run & 8--32 & 18.00 & \raggedright GSM8K accuracy... & \raggedright unset & \raggedright ParScale-1.8B-P8 approx... \tabularnewline
2505.23305 & Run & 8--32 & 10.24 & \raggedright FAD (Fr\'echet Audio... & \raggedright point estimate & \raggedright 0.47 \tabularnewline
2506.01511 & Run & 8--32 & 18.00 & \raggedright Black-box attack success... & \raggedright threshold & \raggedright \textgreater{}50\% black-box ASR (paper... \tabularnewline
2506.18896 & Run & 8--32 & 12.80 & \raggedright MATH500 accuracy & \raggedright threshold & \raggedright \textgreater{}80\% (vs \textasciitilde{}73\%... \tabularnewline
2502.01203 & Retrain & 0--8 & 0.96 & \raggedright Win rate against preferred... & \raggedright direction & \raggedright 66.1\% \tabularnewline
2503.09657 & Retrain & 0--8 & 4.16 & \raggedright Wikitext2 perplexity (lower... & \raggedright threshold & \raggedright 16.17 \tabularnewline
2504.13146 & Retrain & 0--8 & 8.00 & \raggedright distilled student accuracy... & \raggedright point estimate & \raggedright \textasciitilde{}40\% \tabularnewline
2505.15201 & Retrain & 0--8 & 0.00 & \raggedright empirical variance of... & \raggedright magnitude & \raggedright \textasciitilde{}10x lower variance than... \tabularnewline
2505.17836 & Retrain & 0--8 & 0.00 & \raggedright Average normalized absolute... & \raggedright point estimate & \raggedright \textasciitilde{}0.05 \tabularnewline
2505.21077 & Retrain & 0--8 & 1.28 & \raggedright throughput speedup & \raggedright point estimate & \raggedright 1.20x \tabularnewline
2505.24680 & Retrain & 0--8 & 1.28 & \raggedright Average perplexity on... & \raggedright point estimate & \raggedright 24.59 \tabularnewline
2506.06991 & Retrain & 0--8 & 0.00 & \raggedright Area Under ROC Curve (AUC) & \raggedright point estimate & \raggedright 0.91 \tabularnewline
2506.07673 & Retrain & 0--8 & 0.00 & \raggedright estimation gap reduction vs... & \raggedright point estimate & \raggedright 37.2\% \tabularnewline
2506.20233 & Retrain & 0--8 & 0.00 & \raggedright Recovery accuracy... & \raggedright point estimate & \raggedright negligible bias and... \tabularnewline
2509.16391 & Retrain & 0--8 & 0.67 & \raggedright Average Gap (lower is... & \raggedright direction & \raggedright 0.25 \tabularnewline
2510.07501 & Retrain & 0--8 & 0.00 & \raggedright Bias of MR estimator value... & \raggedright point estimate & \raggedright -0.20 \tabularnewline
2510.09485 & Retrain & 0--8 & 0.00 & \raggedright Worst-case total variation... & \raggedright direction & \raggedright Local sampler TV \textless{} Global... \tabularnewline
2510.20725 & Retrain & 0--8 & 3.12 & \raggedright cumulative regret growth... & \raggedright direction & \raggedright sublinear (grows slower... \tabularnewline
2510.21363 & Retrain & 0--8 & 0.17 & \raggedright gender fairness score (0-1... & \raggedright point estimate & \raggedright 0.560 \tabularnewline
2510.23577 & Retrain & 0--8 & 0.64 & \raggedright Average Precision (AP) in... & \raggedright point estimate & \raggedright 98.89 \tabularnewline
2410.19933 & Retrain & 8--32 & 13.00 & \raggedright safety rate (fraction of... & \raggedright threshold & \raggedright \textgreater{}90\% safety rate, above the... \tabularnewline
2505.24089 & Retrain & 8--32 & 12.80 & \raggedright membership-inference AUC on... & \raggedright point estimate & \raggedright \textasciitilde{}82.45\% AUC (paper Table 1... \tabularnewline
2506.08898 & Retrain & 8--32 & 9.60 & \raggedright Normalized Hypervolume (HV) & \raggedright direction & \raggedright 0.6411 \tabularnewline
2506.10943 & Retrain & 8--32 & 12.00 & \raggedright no-context SQuAD QA... & \raggedright direction & \raggedright paper full-scale: base... \tabularnewline
2506.13717 & Retrain & 8--32 & 9.60 & \raggedright Top-1 classification... & \raggedright point estimate & \raggedright 90.21\% \tabularnewline
2510.05874 & Retrain & 8--32 & 15.36 & \raggedright Full Rollout MSE & \raggedright point estimate & \raggedright 1.259e-6 \tabularnewline
2510.10480 & Retrain & 8--32 & 29.44 & \raggedright ISM (Interaction Site... & \raggedright point estimate & \raggedright 52.15\% \tabularnewline
2510.19314 & Retrain & 8--32 & 10.24 & \raggedright Average performance (PERF.) & \raggedright point estimate & \raggedright 0.7923 \tabularnewline
2510.19784 & Retrain & 8--32 & 12.00 & \raggedright test MSE on LV trajectory... & \raggedright threshold & \raggedright paper reports 7.93e-5; bar... \tabularnewline
2510.22123 & Retrain & 8--32 & 10.96 & \raggedright Force MAE (kcal/mol/\AA{}) & \raggedright point estimate & \raggedright 0.102 \tabularnewline
2510.24940 & Retrain & 8--32 & 12.00 & \raggedright answer accuracy on SVAMP... & \raggedright point estimate & \raggedright \textasciitilde{}46\% (paper Table 1... \tabularnewline
2511.00119 & Retrain & 8--32 & 12.00 & \raggedright FID of generated... & \raggedright direction & \raggedright Rectified Flow FID clearly... \tabularnewline
2511.02652 & Retrain & 8--32 & 20.80 & \raggedright SSIM reconstruction quality & \raggedright point estimate & \raggedright 0.8541 \tabularnewline
2504.12463 & Retrain & 32--96 & 35.84 & \raggedright validation perplexity & \raggedright direction & \raggedright 23.48 \tabularnewline
2505.02391 & Retrain & 32--96 & 60.00 & \raggedright Average @ 8 accuracy on... & \raggedright direction & \raggedright \textasciitilde{}71-73\% (vs RAFT++ baseline... \tabularnewline
2506.19839 & Retrain & 32--96 & 48.89 & \raggedright FID 10K & \raggedright point estimate & \raggedright 29.56 \tabularnewline
2508.21046 & Retrain & 32--96 & 36.10 & \raggedright Task success rate & \raggedright point estimate & \raggedright 98.6\% \tabularnewline
2402.04579 & Reimplement & 0--8 & 0.00 & \raggedright Combined cost (modification... & \raggedright direction & \raggedright CCE achieves best combined... \tabularnewline
2410.15392 & Reimplement & 0--8 & 2.00 & \raggedright PSNR (novel view synthesis)... & \raggedright direction & \raggedright up to 3dB higher PSNR and... \tabularnewline
2503.02809 & Reimplement & 0--8 & 0.00 & \raggedright Sharpness S(theta)... & \raggedright point estimate & \raggedright Sharpness oscillates around... \tabularnewline
2505.12677 & Reimplement & 0--8 & 0.60 & \raggedright LPIPS\_e (erasure LPIPS)... & \raggedright threshold & \raggedright 0.41 (Table 1, Kelly... \tabularnewline
2505.19087 & Reimplement & 0--8 & 0.08 & \raggedright Generalization bound value... & \raggedright threshold & \raggedright 0.60302 \tabularnewline
2505.22860 & Reimplement & 0--8 & 1.00 & \raggedright DDI AUC-ROC & \raggedright threshold & \raggedright 0.99 \tabularnewline
2505.24452 & Reimplement & 0--8 & 0.86 & \raggedright validation accuracy (\%) & \raggedright point estimate & \raggedright 94.57$\pm$0.33 \tabularnewline
2506.17475 & Reimplement & 0--8 & 0.10 & \raggedright Test accuracy & \raggedright point estimate & \raggedright 79.01\% \tabularnewline
2507.06489 & Reimplement & 0--8 & 2.56 & \raggedright \% affected samples with... & \raggedright threshold & \raggedright 7.0\% \tabularnewline
2509.16950 & Reimplement & 0--8 & 1.60 & \raggedright Poisoned MVR (mission... & \raggedright point estimate & \raggedright 1.00 \tabularnewline
2510.05767 & Reimplement & 0--8 & 4.60 & \raggedright convergence speed =... & \raggedright direction & \raggedright Greedy-64 converges faster... \tabularnewline
2510.08177 & Reimplement & 0--8 & 0.11 & \raggedright Top-1 accuracy (\%) & \raggedright point estimate & \raggedright 65.9\% \tabularnewline
2510.25146 & Reimplement & 0--8 & 0.05 & \raggedright PSNR (peak signal-to-noise... & \raggedright point estimate & \raggedright 25.8 \tabularnewline
2511.20906 & Reimplement & 0--8 & 2.22 & \raggedright task success rate & \raggedright threshold & \raggedright 100\% \tabularnewline
2512.02339 & Reimplement & 0--8 & 0.12 & \raggedright J\&F mean (Jaccard +... & \raggedright point estimate & \raggedright 77.6\% \tabularnewline
2602.03066 & Reimplement & 0--8 & 4.00 & \raggedright NTK availability A(K,g) =... & \raggedright direction & \raggedright availability(shortcut) \textgreater{}... \tabularnewline
2411.06890 & Reimplement & 8--32 & 23.04 & \raggedright Structural Hamming Distance... & \raggedright direction & \raggedright 1.51 \tabularnewline
2502.08924 & Reimplement & 8--32 & 12.00 & \raggedright test accuracy (greedy... & \raggedright point estimate & \raggedright 0.456 \tabularnewline
2505.10039 & Reimplement & 8--32 & 24.00 & \raggedright circuit completeness =... & \raggedright direction & \raggedright Ns+Dn yields higher... \tabularnewline
2505.16927 & Reimplement & 8--32 & 9.60 & \raggedright IFEval principle-following... & \raggedright threshold & \raggedright 65.6\% \tabularnewline
2506.00070 & Reimplement & 8--32 & 15.36 & \raggedright Spatial reasoning score... & \raggedright point estimate & \raggedright 1.51 \tabularnewline
2507.02064 & Reimplement & 8--32 & 16.00 & \raggedright L2 distance from decoded... & \raggedright point estimate & \raggedright \textasciitilde{}0 cm after convergence... \tabularnewline
2510.25529 & Reimplement & 8--32 & 13.48 & \raggedright Total Average Return (TAR) & \raggedright point estimate & \raggedright 891.7$\pm$19.1 \tabularnewline
2511.01463 & Reimplement & 8--32 & 12.80 & \raggedright R precision Top-3 & \raggedright point estimate & \raggedright 0.785 \tabularnewline
2511.02225 & Reimplement & 8--32 & 18.00 & \raggedright policy learning success rate & \raggedright point estimate & \raggedright 0.95$\pm$0.03 \tabularnewline
2503.14698 & Reimplement & 32--96 & 78.08 & \raggedright PSNR (dB) & \raggedright threshold & \raggedright 30.61 \tabularnewline
2505.22596 & Reimplement & 32--96 & 64.00 & \raggedright gIoU on ReasonSeg zero-shot... & \raggedright point estimate & \raggedright 60.2\% \tabularnewline
2506.02882 & Reimplement & 32--96 & 43.52 & \raggedright Intersection over Union... & \raggedright point estimate & \raggedright 77.0\% \tabularnewline
2506.20024 & Reimplement & 32--96 & 47.36 & \raggedright CRPS (mean across... & \raggedright point estimate & \raggedright 0.904$\times$10$^{-2}$ (0.00904) \tabularnewline
2510.04136 & Reimplement & 32--96 & 56.00 & \raggedright exact & \raggedright direction & \raggedright exact \tabularnewline
2510.15194 & Reimplement & 32--96 & 48.64 & \raggedright Downstream classification... & \raggedright direction & \raggedright 89.68 \tabularnewline
2510.23429 & Reimplement & 32--96 & 96.00 & \raggedright Median Chamfer Distance... & \raggedright point estimate & \raggedright 0.20 \tabularnewline
2511.19808 & Reimplement & 32--96 & 32.40 & \raggedright Test accuracy (\%) & \raggedright point estimate & \raggedright 97.3 \tabularnewline
\end{longtable}
}

% src: outputs/v6/app_rebuild/dev_split.jsonl
{\scriptsize
\begin{longtable}{@{}l@{\hspace{4pt}}l@{\hspace{4pt}}>{\centering\arraybackslash}p{0.9cm}@{\hspace{4pt}}>{\raggedleft\arraybackslash}p{1.1cm}@{\hspace{8pt}}p{3.2cm}@{\hspace{5pt}}p{1.6cm}@{\hspace{5pt}}p{3.3cm}@{}}
\toprule
arXiv & Tier & Band & H100-h & Metric & Bar & Target value \\
\midrule
\endfirsthead
\multicolumn{7}{@{}l}{\footnotesize\itshape Table~\ref{tab:dev14}, continued.} \\[3pt]
\toprule
arXiv & Tier & Band & H100-h & Metric & Bar & Target value \\
\midrule
\endhead
\midrule
\endfoot
\bottomrule
\noalign{\vskip 6pt}
\caption{The 14 development papers, in the same columns, sort order, and truncation as Table~\ref{tab:eval100}. This split is disjoint from the evaluation set and is used for harness development only.}
\label{tab:dev14} \\
\endlastfoot
2505.18513 & Run & 0--8 & 0.01 & \raggedright LDS (Linear Datamodeling... & \raggedright point estimate & \raggedright 21.11 \tabularnewline
2506.09045 & Run & 0--8 & $<$0.01 & \raggedright end-to-end latency speedup & \raggedright point estimate & \raggedright 2.68x \tabularnewline
2507.02546 & Run & 0--8 & 0.01 & \raggedright Metric point map relative... & \raggedright threshold & \raggedright 4.44\% \tabularnewline
2510.21323 & Run & 0--8 & 0.01 & \raggedright Unified concept extraction... & \raggedright point estimate & \raggedright Unified concept set shared... \tabularnewline
2511.07099 & Run & 0--8 & 0.01 & \raggedright Speaker similarity score... & \raggedright point estimate & \raggedright 0.113 \tabularnewline
2502.06067 & Retrain & 0--8 & 2.00 & \raggedright 95\% CI coverage rate & \raggedright threshold & \raggedright 1.0 (100\%) \tabularnewline
2502.08101 & Retrain & 0--8 & 1.00 & \raggedright Mean accuracy (\%) & \raggedright point estimate & \raggedright 94.98\% \tabularnewline
2505.19458 & Retrain & 0--8 & 2.50 & \raggedright Sudoku board accuracy on... & \raggedright point estimate & \raggedright 34.4\% \tabularnewline
2505.13431 & Retrain & 8--32 & 12.00 & \raggedright Success rate (\%) & \raggedright point estimate & \raggedright 98.0\% \tabularnewline
2505.17282 & Reimplement & 0--8 & 2.00 & \raggedright Correlation coefficient... & \raggedright point estimate & \raggedright Clear separation: positive... \tabularnewline
2506.09518 & Reimplement & 0--8 & 2.56 & \raggedright PSNR (dB) & \raggedright point estimate & \raggedright 39.38 \tabularnewline
2510.25739 & Reimplement & 0--8 & 1.80 & \raggedright Inference acceleration... & \raggedright direction & \raggedright 1.71$\times$ \tabularnewline
2511.10107 & Reimplement & 0--8 & 2.00 & \raggedright D1-all error rate (\%) & \raggedright direction & \raggedright 2.77 \tabularnewline
2512.03528 & Reimplement & 0--8 & 1.00 & \raggedright average episode cumulative... & \raggedright point estimate & \raggedright 138.0$\pm$88.1 \tabularnewline
\end{longtable}
}

